\documentclass[12pt]{article}
%^ Start of document
\usepackage{times}  %Makes text TimesNewRoman
\usepackage{setspace} %Allows you to set single or double space 
\usepackage{biblatex} %Imports biblatex package
\usepackage{graphicx} % Required for inserting images
\usepackage{authblk}
\usepackage[colorlinks=true, urlcolor=blue, linkcolor=blue]{hyperref}
\usepackage{subfig}
\usepackage{float}
\usepackage{tabularx}
\usepackage{multirow}
\usepackage[paperheight=11in,
   paperwidth=8.5in,
   top=20mm,
   bottom=20mm,
   left=40mm,
   right=40mm]{geometry}
\usepackage{moresize}
\usepackage{tikz}
\usepackage{xcolor}
\usepackage{physics}
\usepackage{amssymb}
\usepackage{mathrsfs}
\usepackage{amsmath}
\usepackage{ragged2e}
\usepackage{comment}
\usepackage{xcolor}
\addbibresource{classicalHW4.bib}
\begin{document}
\singlespacing  %Sets single spacing for whole document
\title{Classical Mechanics HW 4}

\author[1]{Zachary Tullier}
\affil[1]{Student\\Physics Department\\ University of Houston - Clear Lake \\Houston, Texas\\U.S}

\maketitle

\section{Question 1}
    Find the principal moments of inertia about the center of mass of a flat rigid body in the shape of a $45\deg$ right triangle with uniform mass density. What are the principal axes? Find the spin kinetic energy and the angular momentum.

    \underline{Answer 1, part a}
    The moment of inertia is given by
    \[
        I_{ij} = \int_V \rho \left( \delta_{ij} r^2 - x_i x_j \right) dV
    \]
    Where $r^2 = x^2 + y^2$. I am assuming the inertia is in the body reference frame, so $I_{xy} = I_{xz} = I_{yz} = 0$
    \newline For the x-axis
    \[
        I_{xx} = \int_A y^2 \sigma dA =
    \]
    Expand the area integral into an x and a y integral
    \[
        I_{xx} = \sigma \int_0^a \int_0^{a-x} y^2 dy dx
    \]
    Integrate
    \[
        I_{xx} = \frac{\sigma a^4}{3} (1 - \frac{3}{2} + 1 - \frac{1}{4}) = \frac{\sigma a^4}{12}
    \]
    Apply the definition of density: $\sigma \equiv \frac{2 M}{a^2}$
    \[
        I_{xx} = \frac{\left(\frac{2 M}{a^2} \right) a^4}{12}
    \]
    Simplify
    \[
        I_{xx} = \frac{M a^2}{18}
    \]
    Similarly, for the y-axis
    \[
        I_{yy} = \int_A x^2 \sigma dA =
    \]
    Expand the area integral into an x and a y integral
    \[
        I_{yy} = \sigma \int_0^a \int_0^{a-y} x^2 dy dx
    \]
    Integrate
    \[
        I_{yy} = \frac{\sigma a^4}{3} (1 - \frac{3}{2} + 1 - \frac{1}{4}) = \frac{\sigma a^4}{12}
    \]
    Apply the definition of density: $\sigma \equiv \frac{2 M}{a^2}$
    \[
        I_{yy} = \frac{\left(\frac{2 M}{a^2} \right) a^4}{12}
    \]
    Simplify
    \[
        I_{yy} = \frac{M a^2}{18}
    \]
    Similarly for the z-axis
    \[
        I_{zz} = \int_A z^2 \sigma dA =
    \]
    Expand the area integral into an x and a y integral
    \[
        I_{zz} = \sigma \left( \int_0^a \int_0^{a-y} x^2 dy dx + \int_0^a \int_0^{a-y} y^2 dy dx \right)
    \]
    Integrate
    \[
        I_{zz} = \sigma \left(\frac{a^4}{12} + \frac{3 a^4}{2} \right)
    \]
    Apply the definition of density: $\sigma \equiv \frac{2 M}{a^2}$
    \[
        I_{zz} = \left(\frac{2 M}{a^2} \right) \left(\frac{a^4}{12} + \frac{3 a^4}{2} \right)
    \]
    Simplify
    \[
        I_{zz} = \frac{M a^2}{9}
    \]

    \underline{Answer 1, part b}
    The principal axis run along the base, the height, and perpendicular 

    \underline{Answer 1, part c}
    In the body frame the spin kinetic energy is defined as 
    \[
        T_3 = \frac{1}{2} \omega^{a'} I_{a'b'} \omega^{b'}
    \]
    We can calculate the spin kinetic energy for all three axis (x, y, z).
    \newline The x and y axis are nearly identical
    \[
        T_{3x} = \frac{1}{2} \omega_{x'}^2 I_{xx} = \frac{1}{2} \omega_{x'}^2 \left(\frac{M a^2}{18} \right) = \frac{\omega_{x'}^2 M a^2}{36}
    \]
    \[
        T_{3y} = \frac{1}{2} \omega_{y'}^2 I_{yy} = \frac{1}{2} \omega_{y'}^2 \left(\frac{M a^2}{18} \right) = \frac{\omega_{y'}^2 M a^2}{36}
    \]
    \[
        T_{3z} = \frac{1}{2} \omega_{z'}^2 I_{zz} = \frac{1}{2} \omega_{z'}^2 \left(\frac{M a^2}{9} \right) = \frac{\omega_{z'}^2 M a^2}{18}
    \]
    
    \underline{Answer 1, part d}
    Angular momentum is defined as
    \[
        L = I \omega
    \]
    We can calculate the angular momentum for all three axis (x, y, z)
    \newline The x and y axis are nearly identical
    \[
        L = I_{xx} \omega_x = \frac{M a^2}{18} \omega_x
    \]
    \[
        L = I_{yy} \omega_y = \frac{M a^2}{18} \omega_y
    \]
    \[
        L = I_{zz} \omega_z = \frac{M a^2}{9} \omega_z     
    \]

\section{Question 2}
    The total angular velocity vector in the body frame of a rigid body in terms of the Euler angles and their time derivatives is given by:
    \[
        \omega = \left(\omega^{x'}, \omega^{y'}, \omega^{z'} \right) = \left(\dot{\varphi} sin(\theta) sin(\psi) + \dot{\theta} cos(\psi), \dot{\varphi} sin(\theta) cos(\psi) - \dot{\theta} sin(\psi), \dot{\varphi} cos(\theta) + \dot{\varphi} \right)
    \]
    Show that in the laboratory frame itis given by:
    \[
        \omega = \left(\omega^{x}, \omega^{y}, \omega^{z} \right) = \left(\dot{\psi} sin(\theta) sin(\varphi) + \dot{\theta} cos(\varphi), -\dot{\psi} sin(\theta) cos(\varphi) + \dot{\theta} sin(\varphi), \dot{\psi} cos(\theta) + \dot{\psi} \right)
    \]

    \underline{Answer 2}
    We will use Rotation Matricies to transform the the angular velocity from the body from to the lab frame. We will consider the given set of equations for $\omega$ to be the body frame.
    \[
        \omega_{lab} = R_t \omega_{body}
    \]
    The rotation matricies are defined as:
    \[
        R_t = R_1(\varphi) R_2(\theta) R_3(\psi)
    \]
    \[
        R_1(\varphi) =
        \begin{pmatrix}
            cos(\varphi) & -sin(\varphi) & 0
            \\ sin(\varphi) & cos(\varphi) & 0
            \\ 0 & 0 & 1
        \end{pmatrix}
    \]
    \[
        R_2(\theta) =
        \begin{pmatrix}
            1 & 0 & 0
            \\ 0 & cos(\theta) & sin(\theta)
            \\ 0 -sin(\theta) & cos(\theta)
        \end{pmatrix}
    \]
    \[
        R_3(\psi) =
        \begin{pmatrix}
            cos(\psi) & sin(\psi) & 0
            \\ -sin(\psi) & cos(\psi) & 0
            \\ 0 & 0 & 1
        \end{pmatrix}
    \]
    Multiply  the matricies to get
    \[{\tiny
        R_{tot} =
        \begin{pmatrix}
            cos(\theta) \left(cos(\phi) cos(\psi) - sin(\phi) sin(\phi)\right)
            &
            \left( -sin(\psi)cos(\phi) - cos(\psi) sin(\phi) \right)
            & 
            -sin(\theta) \left(cos(\phi) cos(\psi) + sin(\phi) sin(psi) \right)
            \\
            cos(\theta) \left(sin(\phi) cos(\psi) + cos(\phi) sin(\psi) \right)
            & 
            \left( -sin(\phi) sin(\psi) + cos(\phi)cos(\psi) \right)
            & 
            -sin(\theta) \left( sin(\phi) cos(\psi) + cos(\phi) sin(\psi) \right)
            \\
            sin(\theta)
            &
            0
            &
            cos(\theta)
        \end{pmatrix}
        }
    \]
    Multiply the angular velocity by the rotational matrices for each axis
    \[
        \begin{split}
            \omega_{lab}^x =
            \\
            cos(\theta) \left(cos(\phi) cos(\psi) - sin(\phi) sin(\phi)\right) \left(\dot{\varphi} sin(\theta) sin(\psi) + \dot{\theta} cos(\psi) \right) 
            \\
            + \left( -sin(\psi)cos(\phi) - cos(\psi) sin(\phi) \right) \left(\dot{\varphi} sin(\theta) cos(\psi) - \dot{\theta} sin(\psi) \right)
            \\
            - sin(\theta) \left(cos(\phi) cos(\psi) + sin(\phi) sin(psi) \right) \left(\dot{\varphi} cos(\theta) + \dot{\varphi} \right)
        \end{split}
    \]
    \[
        \begin{split}
            \omega_{lab}^y =
            \\
            cos(\theta) \left(sin(\phi) cos(\psi) + cos(\phi) sin(\psi) \right) \left(\dot{\varphi} sin(\theta) sin(\psi) + \dot{\theta} cos(\psi) \right) 
            \\ 
            + \left( -sin(\phi) sin(\psi) + cos(\phi)cos(\psi) \right) \left(\dot{\varphi} sin(\theta) cos(\psi) - \dot{\theta} sin(\psi) \right)
            \\
            + -sin(\theta) \left( sin(\phi) cos(\psi) + cos(\phi) sin(\psi) \right) \left(\dot{\varphi} cos(\theta) + \dot{\varphi} \right)
        \end{split}
    \]
    \[
        \begin{split}
            \omega_{lab}^z =
            \\
            sin(\theta) \left(\dot{\varphi} sin(\theta) sin(\psi) + \dot{\theta} cos(\psi) \right) 
            \\ 
            + 0
            \\
            + cos(\theta) \left(\dot{\varphi} cos(\theta) + \dot{\varphi} \right)
        \end{split}
    \]
    Simplify to get
    \[
        \omega_{lab}^x = \dot{\psi} sin(\theta) sin(\varphi) + \dot{\theta} cos(\varphi)
    \]
    \[
        \omega_{lab}^y = -\dot{\psi} sin(\theta) cos(\varphi) + \dot{\theta} sin(\varphi)
    \]
    \[
        \omega_{lab}^z = \dot{\psi} cos(\theta) + \dot{\psi}
    \]
    
\newpage
%\printbibliography
\end{document}
